\section{Related Work}\label{s:related}

A number of quantum network simulators have been developed to study entanglement distribution and quantum key distribution. These tools differ in abstraction level, ranging from circuit-based quantum simulators to discrete-event network simulators like this one. 

In the following sections, we discuss SeQUeNCe, Qunetsim, QKDNetsim, SimulaQron, Qiskit, and NetSquid, focusing on their main design choices.

\textbf{Sequence} \cite{wu2021sequence} is the most similar simulator to this work, as it served as the primary inspiration for its design. According to the authors, its main goal is to provide a framework suitable for simulating quantum network prototypes for both current and future technologies.

Although our simulator differs in scope, as this one is specifically designed for the E91 QKD protocol, the two systems share several architectural similarities. \textbf{Sequence} is an event-driven simulator that supports precise event scheduling with timestamp accuracy. It also follows a layered architecture with a clear separation of concerns.

However, \textbf{Sequence} extends this design with additional abstraction layers, such as dedicated network and photon layers. Furthermore, its fidelity modelling differs from ours: instead of applying a constant degradation factor per kilometer, it incorporates a combination of quantum memory fidelity degradation and additional physical parameters to more accurately represent realistic quantum network behavior.

Another example of a layered simulator is QKDNetSim \cite{mehic2022qkdnetsim}. This simulator provides multiple layers, such as the key management layer and the quantum layer. It was introduced in 2017 as the first simulation environment for quantum key distribution networks. The main similarity with this simulator is that its primary focus is on quantum key distribution.

However, QKDNetSim differs in that it does not model quantum physical behavior. Instead, it simulates only the packet-based communication required for QKD between nodes. For this reason, QKDNetSim is built on top of ns-3, a discrete-event, classical network simulator.

\textbf{QunetSim} \cite{diadamo2021qunetsim},similar to SeQUeNCe, it provides a flexible, event-driven framework in which users define network topologies and protocols programmatically rather than through a dedicated user interface. QunetSim also supports larger quantum networks through intermediary hosts and entanglement swapping, enabling the simulation of multi-hop entanglement distribution. Unlike this work, however, it is designed as a general-purpose quantum networking framework rather than a simulator focused specifically on the E91 QKD protocol.
